\scanpage{26}%
% archive sheet number: 25
tuelles) permet de définir des sections $a_i,a'_i$ de $V,\check V$ -- en fait, $a'_0$ est une section de $L'_0$ définie en termes de \deleted{$D_{s_0}$} la \deleted{droite} \inserted{section} $s_0$ de $\unsure{E|\mathcal E}$ \deleted{$\Delta_{\unsure{s|\xi}}\subset E$ \unsure{engendrée par} $s_0$ et} \inserted{(au dessus de $1$,} définissant un \unsure{splittage|splitting} de $E$ en $E\simeq V\oplus\underline{\mathcal O}_S$ \ldots) par $\langle s_0,a'_0\rangle=1$, ce qui permet ensuite de définir de proche en proche les bases $a_0,a'_1,a_{\unsure{2|1}},\ldots,a'_n,a_n$ de $L_0,L'_1,L_1,\ldots$ par les conditions
\[
\tau_i^E=\mathrm{id}+a'_i\otimes a_i,\quad
\langle a_i,a'_{i+1}\rangle=1\qquad 0\leq i\leq n
\]
(\uline{NB} il y a lieu de poser $a_{-1}=s_0$\unsure{.})

Inversement, donnons nous une base $a'_0$ de $L'_0$. \deleted{Considérons \unsure{ce que veut|indique voir}} le plan
\[
F=\bigcap_{1\leq i\leq n}\mathcal H_i\qquad
\text{orthogonal aux}\quad L'_i\quad 1\leq i\leq n,
\]
\deleted{dans celui-ci l'équation} on a \deleted{$\mathcal H_0\cap F$} $F\not\subset\mathcal H_0$ en tte fibre i.e. $F\cap\mathcal H_0$ est une droite (la droite $\Delta_{\unsure{o|0|c}}$ \deleted{\unsure{engendrée}} $\subset E$ correspondant à $o$), et l'équation \deleted{$\unsure{\xi}$} $\langle s,a'_0\rangle=1$ définit dans $F$ une \deleted{\unsure{translat.}} \inserted{droite affine} \deleted{\unsure{lin.}} translatée de $\Delta$, soit $\Delta'$ (dépendant du choix de $a'_0$ -- \deleted{\unsure{il y a une bijection}} (\uline{NB} : $L'_0\simeq F/\Delta$)

Considérons $F\xrightarrow{\pi|F}\underline{\mathcal O}_S$
